\scanpage{63}
% archive sheet number: 62
\uline{Groupe de Brauer global.}

Pour avoir des fondements solides et bien connus, nous nous plaçons en Géométrie \deleted{Algébrique} \inserted{analytique} [sur le corps $\mathbb{C}$ des complexes] Les groupes de cohomologie seront pris p.r. à la topologie usuelle.

Du \deleted{\unsure{les}} \inserted{diagramme de} suites exactes
\diagram{\begin{tikzcd}[ampersand replacement=\&,column sep=small]
e\arrow[r] \& \mathbb{C}^*\arrow[r] \& \operatorname{Gl}(n)\arrow[r] \& \operatorname{GP}(n-1)\arrow[r] \& e \\
e\arrow[r] \& \mathbb{Z}/n\mathbb{Z}\arrow[u]\arrow[r,dash] \& \operatorname{Sl}(n)\arrow[u]\arrow[r] \& \deleted{\unsure{\operatorname{Gl}}}\operatorname{GP}(n-1)\arrow[u]\arrow[r,dash] \& e
\end{tikzcd}}
on tire (pour un \unsure{bon} $X$) la ``suite exacte''
\diagram{\begin{tikzcd}[ampersand replacement=\&,column sep=tiny]
\cdots\arrow[r] \& H^1(X,\underline{\mathbb{C}}^*)\arrow[r] \& H^1(X,\underline{\operatorname{Gl}}(n))\arrow[r] \& H^1(X,\operatorname{GP}(n-1))\arrow[r] \& H^2(X,\mathbb{C}^*) \\
\cdots\arrow[r] \& H^1(X,\mathbb{Z}/n\mathbb{Z})\arrow[u]\arrow[r] \& H^1(X,\underline{\operatorname{Sl}}(n))\arrow[u]\arrow[r] \& H^1(X,\operatorname{GP}(n-1))\arrow[u]\arrow[r,dash] \& H^2(X,\mathbb{Z}/n\mathbb{Z})\arrow[u]
\end{tikzcd}}
Considérons, pour un \unsure{fibré} donné $\unsure{\omega|\sigma|\varpi}\in H^1(X,\operatorname{GP}(n-1))$, l'élément
\[
\delta(\omega)\in H^2(X,\mathbb{Z}/n\mathbb{Z})
\]
obstruction à remonter le groupe structural à $\operatorname{Sl}(n)$. Alors l'obstruction à remonter le groupe structural à $\operatorname{Gl}(n)$ est l'image de $\delta(\omega)$ dans $H^2(X,\mathbb{C}^*)$; elle est donc d'ordre $n$ [l'ensemble des élémts d'ordre $n$ dans $H^2(X,\mathbb{C}^*)$ étant exactement l'image de $H^2(X,\mathbb{Z}/n\mathbb{Z})$].

Les \unsure{classes} de fibrés sur $X$, de groupe $\operatorname{GP}(n-1)$, sont aussi les \unsure{classes} de fibrés en algèbres \deleted{\unsure{matrices}} de \deleted{\unsure{dim.}} \inserted{\unsure{r.g|r}} $n^2$; celles qui proviennent de $\operatorname{Gl}(n)$ sont les fibrés en algèbres qui
